% ==================================================
% CHAPTER 3: SYSTEM DESIGN
% ==================================================

\chapter{System Design}
\label{ch:design}

\section{Architecture Overview}

The service is a layered Python application behind a single-page web
interface (Figure~\ref{fig:arch}). The layers have strict rules about
what they may depend on:

\begin{itemize}
  \item The \textbf{domain engine} (\texttt{hostelswap/domain}) holds all
  the matching mathematics and the proposal state machine. It imports no
  web framework and no database driver, so it can be tested with no
  infrastructure at all.
  \item The \textbf{data-access layer} (\texttt{db.py}) is the only
  module that contains SQL. It loads and stores the same data shapes the
  engine uses.
  \item The \textbf{workflow layer} (\texttt{rounds.py}) wires the two
  together and enforces the round rules (``priorities cannot change once
  the round is locked''). It contains no SQL and no matching logic.
  \item The \textbf{API layer} (\texttt{api.py}, \texttt{auth.py})
  authenticates requests, checks roles and translates between JSON and
  the workflow. It also serves the compiled React interface, so one
  process serves the whole application.
\end{itemize}

\begin{figure}[H]
\centering
\resizebox{\textwidth}{!}{\begin{tikzpicture}
  % ---- browser -------------------------------------------------------------
  \node[comp, minimum width=32mm] (stu) at (0, 1.1) {Student pages\\\scriptsize priorities, roommate,\\\scriptsize options, proposals};
  \node[comp, minimum width=32mm] (adm) at (0,-1.1) {Warden pages\\\scriptsize room types, rounds,\\\scriptsize summary, proposals};
  \begin{scope}[on background layer]
    \node[group, fill=blue!3, fit=(stu)(adm), inner sep=7pt] (browser) {};
  \end{scope}
  \node[font=\small\bfseries, anchor=south] at (browser.north) {Browser (React)};

  % ---- server --------------------------------------------------------------
  \node[comp, minimum width=27mm, minimum height=38mm] (api) at (4.3, 0)
    {\textbf{FastAPI}\\\texttt{api.py}\\[1mm]\scriptsize routes, roles\\\scriptsize \texttt{auth.py}\\\scriptsize sessions\\\scriptsize static UI build};
  \node[comp, minimum width=25mm, minimum height=38mm] (rounds) at (7.5, 0)
    {\textbf{Workflow}\\\texttt{rounds.py}\\[1mm]\scriptsize round states\\\scriptsize offer, respond\\\scriptsize expiry sweep};
  \node[comp, minimum width=30mm] (dom) at (11.1, 1.25)
    {\textbf{Domain engine}\\\scriptsize\texttt{domain/}: scoring,\\\scriptsize Hungarian, Murty,\\\scriptsize cycles, Pareto, ranking};
  \node[comp, minimum width=30mm] (dbm) at (11.1,-1.25)
    {\textbf{Data access}\\\scriptsize\texttt{db.py} (only SQL)\\\scriptsize connection pool};
  \begin{scope}[on background layer]
    \node[group, fill=gray!3, fit=(api)(rounds)(dom)(dbm), inner sep=7pt] (server) {};
  \end{scope}
  \node[font=\small\bfseries, anchor=south] at (server.north) {Python service (\texttt{hostelswap})};

  \node[db, minimum width=34mm] (pg) at (11.1,-3.9) {PostgreSQL (Supabase)\\\scriptsize 15 tables, triggers};
  \node[ext] (demo) at (4.3,-3.9) {Demo cohort\\\scriptsize\texttt{demo\_data.py}};

  % ---- flows ---------------------------------------------------------------
  \draw[flow] (stu.east) -- node[lbl, above] {JSON} (api.west |- stu.east);
  \draw[flow] (adm.east) -- node[lbl, above] {JSON} (api.west |- adm.east);
  \draw[flow] (api) -- (rounds);
  \draw[flow] (rounds.east |- dom) -- (dom.west);
  \draw[flow] (rounds.east |- dbm) -- (dbm.west);
  \draw[flow, <->] (dbm) -- node[lbl, right] {SQL} (pg);
  \draw[flow, dashed] (api) -- node[lbl, right] {no database} (demo);
\end{tikzpicture}
}
\caption{Architecture of the Hostel Room Swapping Service}
\label{fig:arch}
\end{figure}

When no database is configured, the API falls back to a fixed 30-student
demo cohort, so the engine can be tried and tested on any machine. In
production this fallback is refused (Section~\ref{sec:failfast}).

\section{Data Flow}

\subsection{Context Diagram (Level 0)}

Figure~\ref{fig:dfd0} shows the system and its two external entities.

\begin{figure}[H]
\centering
\resizebox{0.95\textwidth}{!}{\begin{tikzpicture}
  \node[proc, minimum size=30mm] (sys) at (0,0) {\textbf{0}\\ Hostel Room\\ Swapping\\ Service};
  \node[box, fill=gray!15, minimum width=24mm] (stu) at (-6,0) {Student};
  \node[box, fill=gray!15, minimum width=24mm] (war) at (6,0) {Warden};
  \draw[flow] ([yshift=4mm]stu.east) -- node[lbl, above, align=center] {priorities, roommate\\ requests, offers, replies} ([yshift=4mm]sys.west);
  \draw[flow] ([yshift=-4mm]sys.west) -- node[lbl, below, align=center] {swap options,\\ proposal status} ([yshift=-4mm]stu.east);
  \draw[flow] ([yshift=4mm]war.west) -- node[lbl, above, align=center] {room types,\\ round commands} ([yshift=4mm]sys.east);
  \draw[flow] ([yshift=-4mm]sys.east) -- node[lbl, below, align=center] {round summary,\\ proposals} ([yshift=-4mm]war.west);
\end{tikzpicture}
}
\caption{Context diagram (DFD level 0)}
\label{fig:dfd0}
\end{figure}

\subsection{Level 1}

Figure~\ref{fig:dfd1} splits the system into its four processes.
Registration writes preferences (D1); matching reads them and writes
candidate chains (D2); students offer and answer chains, which become
proposals (D3); and an approved proposal appends new allocation rows
(D4).

\begin{figure}[H]
\centering
\resizebox{0.95\textwidth}{!}{\begin{tikzpicture}
  \node[box, fill=gray!15, minimum width=22mm] (stu) at (-6.6, 0.6) {Student};
  \node[box, fill=gray!15, minimum width=22mm] (war) at ( 6.6, 0.6) {Warden};

  \node[proc] (p1) at (-3.3, 2.6) {\textbf{1}\\ Register\\ preferences};
  \node[proc] (p2) at ( 3.3, 2.6) {\textbf{2}\\ Run\\ matching};
  \node[proc] (p3) at (-3.3,-1.6) {\textbf{3}\\ Offer and\\ approve};
  \node[proc] (p4) at ( 3.3,-1.6) {\textbf{4}\\ Execute\\ swap};

  \node[store, minimum width=40mm] (d1) at (0, 4.4)  {D1\hspace{5mm} enrolments, preferences};
  \node[store, minimum width=40mm] (d2) at (0, 0.6)  {D2\hspace{5mm} chain options};
  \node[store, minimum width=40mm] (d3) at (0,-3.4)  {D3\hspace{5mm} proposals, members};
  \node[store, minimum width=40mm] (d4) at (6.6,-3.4) {D4\hspace{5mm} allocations};

  \draw[flow] (stu) -- node[lbl, left] {priorities} (p1);
  \draw[flow] (p1) |- (d1.west);
  \draw[flow] (d1.east) -| (p2);
  \draw[flow] (war) -- node[lbl, right] {lock and run} (p2);
  \draw[flow] (p2) |- node[lbl, pos=0.75, above] {chains} (d2.east);
  \draw[flow] (d2.west) -| node[lbl, pos=0.25, above] {my options} (p3);
  \draw[flow] (stu) -- node[lbl, left] {offer, reply} (p3);
  \draw[flow, <->] (p3) |- (d3.west);
  \draw[flow] (d3.east) -- node[lbl, above] {approved} (p4);
  \draw[flow] (p4) -- node[lbl, right] {new rows} (d4);
\end{tikzpicture}
}
\caption{Data flow diagram, level 1}
\label{fig:dfd1}
\end{figure}

\section{Data Model}

\subsection{Rooms, Slots and Allocations}
\label{sec:bed-slot}

A student's preferences are properties of a \emph{room}: hostel, floor,
direction, washroom type and room type. The Hungarian algorithm,
however, needs a one-to-one problem. A three-bed room is not one
assignable unit but three. The schema therefore splits every room into
\textbf{bed slots}, one per bed, and students are assigned to slots. The
assignment stays a clean permutation (as many slots as participating
students) while shared rooms and roommate requests remain expressible.

Who holds which slot is never stored as a mutable field. The
\texttt{allocations} table is an append-only history: each row says that
a student holds a slot from a given time. The current allocation is a
view over that history, and two triggers raise an error on any
\texttt{UPDATE} or \texttt{DELETE}, so a bug in any later client cannot
rewrite the record.

\subsection{Tables}

The schema has 15 tables, created by migration \texttt{0001} and changed
by \texttt{0002} (Table~\ref{tab:tables}).

\begin{table}[H]
\centering
\small
\begin{tabularx}{\textwidth}{>{\ttfamily\raggedright\arraybackslash}p{0.33\textwidth} L}
\toprule
\textrm{\textbf{Table}} & \textbf{Holds} \\
\midrule
students, sessions & Accounts with role and PBKDF2 hash; server-side session tokens \\
hostels, rooms, bed\_slots & Physical layout: floor, direction, washroom, room type, one slot per bed \\
hostel\_room\_type\_inventory & Room types a hostel declares, used for valid priorities \\
allocations & Append-only history of who held which slot from when \\
swap\_rounds, round\_enrollments & A round per hostel and its lifecycle; who joined it \\
preference\_sets, preferences & Weighted hard/soft priorities, one active set per student per round \\
roommate\_requests & Roommate requests and whether they were accepted \\
round\_chain\_options & Candidate chains from one engine run, kept so students can browse them later \\
swap\_proposals, swap\_chain\_members & Offered chains, their state, deadline and each member's reply \\
\bottomrule
\end{tabularx}
\caption{Database tables}
\label{tab:tables}
\end{table}

\subsection{Invariants Enforced by the Database}

Three rules are too important to leave to application code:

\begin{enumerate}
  \item \textbf{Append-only history}: triggers on \texttt{allocations}.
  \item \textbf{One active preference set} per student per round: a
  partial unique index.
  \item \textbf{One live proposal} per student: a partial unique index on
  \texttt{swap\_chain\_members} over rows that have not settled. Without
  it, a single student could sit in several chains, and their one
  refusal could invalidate all of them at once.
\end{enumerate}

The third index originally treated every member row that was not
\texttt{rejected} as live. Section~\ref{sec:expiry-defect} explains why
that was wrong and how migration \texttt{0002} fixed it.

\section{State Machines}

\subsection{Round Lifecycle}

A round for one hostel moves through the states in
Figure~\ref{fig:round-states}. Locking is one-way: it closes
registration and runs the engine in the same request, so priorities can
never change after the engine has read them.

\begin{figure}[H]
\centering
\begin{tikzpicture}[node distance=6mm and 9mm]
  \node[state] (draft) {draft};
  \node[state, right=of draft] (open) {open};
  \node[state, right=of open] (locked) {locked};
  \node[state, right=of locked] (running) {running};
  \node[final, right=of running] (done) {completed};
  \node[final, below=10mm of open] (cancel) {cancelled};
  \draw[flow] (draft) -- node[lbl, above] {open} (open);
  \draw[flow] (open) -- node[lbl, above] {lock} (locked);
  \draw[flow] (locked) -- (running);
  \draw[flow] (running) -- node[lbl, above] {chains} (done);
  \draw[flow, dashed] (draft) |- (cancel);
  \draw[flow, dashed] (open) -- (cancel);
  \draw[flow, dashed] (done) |- (cancel);
\end{tikzpicture}
\caption{Round lifecycle (dashed: a warden cancels; open offers are withdrawn)}
\label{fig:round-states}
\end{figure}

\subsection{Proposal Lifecycle}

One proposal is one chain. Each member's reply is recorded separately,
and the proposal moves as in Figure~\ref{fig:proposal-states}:

\begin{itemize}
  \item the \textbf{first rejection} rejects the whole chain, because a
  cycle only works whole, and frees every member;
  \item the \textbf{last approval} approves the chain, which is executed
  at once and appends one allocation row per member;
  \item a proposal still collecting approvals after its
  \textbf{three-day deadline} expires, and its members are freed.
\end{itemize}

\begin{figure}[H]
\centering
\begin{tikzpicture}[node distance=8mm and 14mm]
  \node[state] (prop) {proposed};
  \node[state, right=of prop] (appr) {approved};
  \node[final, right=of appr] (exec) {executed};
  \node[final, below left=12mm and -4mm of prop] (rej) {rejected};
  \node[final, below right=12mm and -4mm of prop] (exp) {expired};
  \node[final, above=8mm of appr] (canc) {cancelled};
  \draw[flow] (prop) to[loop above] node[lbl, above] {approve (not last)} (prop);
  \draw[flow] (prop) -- node[lbl, above] {last approve} (appr);
  \draw[flow] (appr) -- node[lbl, above] {execute} (exec);
  \draw[flow] (prop) -- node[lbl, left] {any reject} (rej);
  \draw[flow] (prop) -- node[lbl, right] {deadline passed} (exp);
  \draw[flow, dashed] (prop) |- (canc);
  \draw[flow, dashed] (appr) -- (canc);
\end{tikzpicture}
\caption{Proposal lifecycle (dashed: a warden withdraws the ring)}
\label{fig:proposal-states}
\end{figure}

Every transition in the domain returns a \emph{new} proposal object and
mutates nothing. Executing a chain emits fresh allocation rows rather
than editing old ones, which matches the append-only history.

Chains offered from one assignment are disjoint cycles, so they execute
independently. A student dropping out of one chain cannot invalidate any
other chain. This is the answer to the \emph{cycle invalidation}
question: a refusal or a silence is contained in the one chain it
belongs to, and the deadline bounds how long a silence can last.

\section{Interface Design}

\subsection{REST API}

The backend exposes 34 routes: four that serve the interface and 30 JSON
endpoints (Table~\ref{tab:api}). Every student and admin endpoint checks
the session's role and answers \texttt{401} or \texttt{403} otherwise.

\begin{table}[H]
\centering
\small
\begin{tabularx}{\textwidth}{>{\raggedright\arraybackslash}p{0.15\textwidth} >{\ttfamily\scriptsize\raggedright\arraybackslash}p{0.5\textwidth} L}
\toprule
\textbf{Group} & \textrm{\textbf{Example endpoints}} & \textbf{Purpose} \\
\midrule
Auth (3) & POST /api/auth/login, POST /api/auth/logout, GET /api/auth/me & Sign in and out, current user \\
Student (14) & PUT /api/student/round/\{id\}/priorities \newline GET /api/student/round/\{id\}/results \newline POST /api/student/round/\{id\}/chains/offer \newline POST /api/student/proposals/\{id\}/respond & Enrol, priorities, roommate requests, options, offers, replies \\
Admin (11) & PUT /api/admin/hostels/\{id\}/room-types \newline POST /api/admin/rounds/\{id\}/lock-and-run \newline POST /api/admin/rounds/\{id\}/cancel \newline POST /api/admin/proposals/\{id\}/cancel & Room types, round lifecycle and cancellation, summary, proposals \\
Diagnostic (2) & GET /api/schema, GET /api/demo & Live schema check; reproducible demo run \\
\bottomrule
\end{tabularx}
\caption{REST API}
\label{tab:api}
\end{table}

\subsection{User Interface}

The interface is a React and TypeScript single-page application built
with Vite. It has three pages and four dialogs:

\begin{itemize}
  \item \textbf{Login}: one form for both roles; the server decides which
  page the user lands on.
  \item \textbf{Student}: current room, round status, enrolment,
  priorities (\texttt{PreferencesModal}), roommate requests, swap options
  shown as chain cards with match percentage and met and missed
  preferences, offering a chain (\texttt{OfferModal}), and pending
  proposals to accept or reject.
  \item \textbf{Admin}: hostels, declared room types
  (\texttt{RoomTypeModal}), the current round with its lifecycle buttons,
  lock-and-run with an optional chain cap (\texttt{RunRoundModal}), the
  round summary and the proposals.
\end{itemize}

The compiled build is committed to the repository and served by FastAPI,
so a deployment does not need Node.js. During development, the Vite dev
server proxies \texttt{/api} to the backend.

\begin{figure}[H]
\centering
\screenshot[0.9\textwidth]{figures/student.png}{student page with swap options (figures/student.png)}
\caption{Student page}
\label{fig:student}
\end{figure}

\begin{figure}[H]
\centering
\screenshot[0.9\textwidth]{figures/admin.png}{warden page with a round and its proposals (figures/admin.png)}
\caption{Warden page}
\label{fig:admin}
\end{figure}

\section{Design Decisions}
\label{sec:decisions}

Table~\ref{tab:decisions} lists the main design decisions, the
alternatives considered and the reason for each choice.

\begin{longtable}{>{\raggedright\arraybackslash}p{0.22\textwidth} >{\raggedright\arraybackslash}p{0.25\textwidth} >{\raggedright\arraybackslash}p{0.43\textwidth}}
\toprule
\textbf{Decision} & \textbf{Alternatives} & \textbf{Reason} \\
\midrule
\endhead
Closed-pool assignment with chains as cycles & Search for chains directly in a graph & Every permutation splits uniquely into disjoint cycles, so chains come for free and are independent \\
Hungarian algorithm \cite{kuhn1955hungarian} & Greedy matching & Optimal in $O(n^3)$; available in SciPy \cite{virtanen2020scipy} \\
Murty's $k$-best \cite{murty1968ranking} & Hungarian only & The optimum is often rejected by the Pareto filter or needs very long chains; alternatives find executable ones \\
Not Gale--Shapley \cite{gale1962college} & Stable matching & Needs two-sided preferences; rooms do not rank students \\
Not Top Trading Cycles \cite{shapley1974cores} & TTC, as in kidney exchange \cite{roth2004kidney} & Needs strict ranked lists; cannot express weighted preferences over several attributes \\
Pareto filter per chain & Maximise total satisfaction only & A student will not approve a swap that makes them worse off \\
Up to three chains per student & One plan for the hostel & One plan gives most students nothing; per-student options reach many more \\
Chain cap of eight & No cap, or five, or ten & Benchmark trade-off of reach against approvals needed (Chapter~\ref{ch:evaluation}) \\
Three-day proposal deadline & No deadline & A silent member must not lock everyone else out \\
Invariants in the database & Checks in Python only & Triggers and indexes hold for every client, including future ones \\
Session cookies, no external identity provider & JWT, OAuth & Revocable on logout; standard library only; httpOnly cookies are not readable by scripts \\
Numbered migrations & Edit one \texttt{schema.sql} & Every database reaches the same schema in the same order; applied files never change \\
\bottomrule
\caption{Design decisions}
\label{tab:decisions}
\end{longtable}
